\chapter{Literatur-, Abbildungs- und Tabellenverzeichnis}\label{literatur--abbildungs--und-tabellenverzeichnis}

\section{Literatur}\label{literatur}

Ackermann, Jan, Katharina Behne, Felix Buchta, Marc Drobot \& Philipp Knopp. 2015. \emph{Metamorphosen des Extremismusbegriffes: Diskursanalytische Untersuchungen zur Dynamik einer funktionalen Unzulänglichkeit}. Wiesbaden: Springer.

Almela-Sánchez, Moisés \& Pascual Cantos-Gómez (eds.). 2018. \emph{Lexical Collocation Analysis: Advances and Applications} (Quantitative Methods in the Humanities and Social Sciences). 1st ed. 2018. Cham: Springer. https://doi.org/10.1007/978-3-319-92582-0.

Altmayer, Claus. 2006. Landeskunde als Kulturwissenschaft: Ein Forschungsprogramm. In Andrea Bogner, Konrad Ehlich, Ludwig M. Eichinger, Andreas F. Kelletat, Hans-Jürgen Krumm, Willy Michel \& Alois Wierlacher (eds.), \emph{Jahrbuch Deutsch als Fremdsprache 32}, 181--199. München: Iudicium.

Anwar, Saba, Artem Shelmanov, Alexander Panchenko \& Chris Biemann. 2020. Generating Lexical Representations of Frames Using Lexical Substitution. In Christine Howes, Stergios Chatzikyriakidis, Adam Ek \& Vidya Somashekarappa (eds.), \emph{Proceedings of the Probability and Meaning Conference (PaM 2020)}, 95--103. Gothenburg: Association for Computational Linguistics. https://aclanthology.org/2020.pam-1.13. (26 March, 2024).

Anwar, Saba, Dmitry Ustalov, Nikolay Arefyev, Simone Paolo Ponzetto, Chris Biemann \& Alexander Panchenko. 2019. HHMM at SemEval-2019 Task 2: Unsupervised Frame Induction using Contextualized Word Embeddings. In Jonathan May, Ekaterina Shutova, Aurelie Herbelot, Xiaodan Zhu, Marianna Apidianaki \& Saif M. Mohammad (eds.), \emph{Proceedings of the 13th International Workshop on Semantic Evaluation}, 125--129. Minneapolis, Minnesota, USA: Association for Computational Linguistics. https://doi.org/10.18653/v1/S19-2018.

Backes, Uwe. 1989. \emph{Politischer Extremismus in demokratischen Verfassungsstaaten}. Wiesbaden: VS Verlag für Sozialwissenschaften. https://doi.org/10.1007/978-3-322-86110-8.

Backes, Uwe. 2018. Extremistische Ideologien. In Eckhard Jesse \& Tom Mannewitz (eds.), \emph{Extremismusforschung: Handbuch für Wissenschaft und Praxis}, 99--159. 1. Auflage. Baden-Baden: Nomos.

Backes, Uwe \& Eckhard Jesse. 1993. \emph{Politischer Extremismus in der Bundesrepublik Deutschland} (Studien Zur Geschichte Und Politik). 3., völlig überarbeitete und aktualisierte Aufl. Vol. 272. Bonn: Bundeszentrale für Politische Bildung.

Backes, Uwe \& Eckhard Jesse. 2005. \emph{Vergleichende Extremismusforschung} (Extremismus und Demokratie). 1. Aufl. Vol. 11. Baden-Baden: Nomos.

Baker, Paul. 2010. \emph{Sociolinguistics and Corpus Linguistics} (Edinburgh Sociolinguistics). Edinburgh: Edinburgh University Press.

Baker, Paul. 2012. Acceptable bias? Using corpus linguistics methods with critical discourse analysis. \emph{Critical Discourse Studies} 9(3). 247--256. https://doi.org/10.1080/17405904.2012.688297.

Baker, Paul. 2016. The Shapes of Collocation. \emph{International Journal of Corpus Linguistics}. John Benjamins Publishing Company 21(2). 139--164. https://doi.org/10.1075/ijcl.21.2.01bak.

Baker, Paul, Gavin Brookes, Dimitrinka Atanasova \& Stuart W. Flint. 2020. Changing frames of obesity in the UK press 2008--2017. \emph{Social Science \& Medicine} 264. https://doi.org/10.1016/j.socscimed.2020.113403.

Baker, Paul, Costas Gabrielatos \& Tony McEnery. 2013. \emph{Discourse Analysis and Media Attitudes: The Representation of Islam in the British Press}. Cambridge: Cambridge University Press. https://doi.org/10.1017/CBO9780511920103.

Barsalou, Lawrence W. 1993. Flexibility, Structure and Linguistic Vagary in Concepts: Manifestations of a Compositional System of Perceptual Symbols. In Alan F. Collins, Susan E. Gathercole, Martin A. Conway \& Peter E. Morris (eds.), \emph{Theories of Memory}, 29--101. Hove, UK; Hillsdale, USA: Erlbaum.

Belica, Cyril. 2001. Kookkurrenzdatenbank CCDB: Eine korpuslinguistische Denk- und Experimentierplattform für die Erforschung und theoretische Begründung von systemisch-strukturellen Eigenschaften von Kohäsionsrelationen zwischen den Konstituenten des Sprachgebrauchs. Mannheim: Institut für Deutsche Sprache.

Belica, Cyril. 2011. Semantische Nähe als Ähnlichkeit von Kookkurrenzprofilen. In Andrea Abel \& Renata Zanin (eds.), \emph{Korpora in Lehre und Forschung}, 155--178. Bozen University Press. https://ids-pub.bsz-bw.de/frontdoor/index/index/docId/2836.

Belica, Cyril \& Kathrin Steyer. 2008. Korpusanalytische Zugänge zu sprachlichem Usus. In \emph{Beiträge zur bilingualen Lexikographie}, 7--24. Univ. Karlova, Filozifická Fak. https://ids-pub.bsz-bw.de/frontdoor/index/index/docId/4265.

Berlin-Brandenburgische Akademie der Wissenschaften (ed.). o.J. DWDS -- Digitales Wörterbuch der deutschen Sprache. https://www.dwds.de/. (5 November, 2019).

Biemann, Chris, Gerhard Heyer \& Uwe Quasthoff. 2022. \emph{Wissensrohstoff Text: Eine Einführung in das Text Mining}. Wiesbaden: Springer Fachmedien Wiesbaden. https://doi.org/10.1007/978-3-658-35969-0.

Bisiada, Mario, Oliver Czulo \& Eleonore Schmitt. 2023. \#MeToo in drei Sprachen: Qualitative Analyse von Konzepten und Diskursmustern im Englischen, Deutschen und Spanischen anhand von Twitter. \emph{Deutsche Sprache} (1). 5. https://doi.org/10.37307/j.1868-775X.2023.01.05.

Blaette, Andreas. 2020a. Compute Log-likelihood Statistics. In \emph{polmineR: Verbs and Nouns for Corpus Analysis. R package version v0.8.5.} https://polmine.github.io/polmineR/reference/ll.html.

Blaette, Andreas. 2020b. GermaParl Corpus of Plenary Protocols. Zenodo. https://doi.org/10.5281/ZENODO.3742113.

Blaette, Andreas \& Christoph Leonhardt. 2020. polmineR: Verbs and Nouns for Corpus Analysis. https://CRAN.R-project.org/package=polmineR. (3 July, 2020).

Blei, David M., Andrew Y. Ng \& Michael I. Jordan. 2003. Latent Dirichlet Allocation. \emph{Journal of Machine Learning Research} (3). 993--1022.

Bötticher, Astrid. 2017. \emph{Radikalismus und Extremismus: Konzeptualisierung und Differenzierung zweier umstrittener Begriffe in der deutschen Diskussion}. Leiden. https://openaccess.leidenuniv.nl/bitstream/handle/1887/49257.

Brezina, Vaclav, Tony McEnery \& Stephen Wattam. 2015. Collocations in Context: A New Perspective on Collocation Networks. \emph{International Journal of Corpus Linguistics} 139--173. https://doi.org/10.1075/ijcl.20.2.01bre.

Brezina, Vaclav \& William Platt. 2023. \#LancsBox X. http://lancsbox.lancs.ac.uk.

Brezina, Vaclav, Pierre Weill-Tessier \& Tony McEnery. 2020. \#LancsBox v.5.x. http://corpora.lancs.ac.uk/lancsbox.

Brookes, Gavin \& Tony McEnery. 2019. The Utility of Topic Modelling for Discourse Studies: A critical Evaluation. \emph{Discourse Studies} 21(1). 3--21. https://doi.org/10.1177/1461445618814032.

Brunner, Annelen \& Kathrin Steyer. 2007. Corpus-Driven Study of Multi-Word Expressions Based on Collocations from a Very Large Corpus. In \emph{Proceedings of the 4th Corpus Linguistics conference, Birmingham}. University of Birmingham. https://ids-pub.bsz-bw.de/frontdoor/index/index/docId/4141.

Brunner, Annelen, Ngoc Duyen Tanja Tu, Lukas Weimer \& Fotis Jannidis. 2020. To BERT or not to BERT -- Comparing Contextual Embeddings in a Deep Learning Architecture for the Automatic Recognition of Four Types of Speech, Thought and Writing Representation. In \emph{Proceedings of the 5th Swiss Text Analytics Conference (SwissText) \& 16th Conference on Natural Language Processing (KONVENS)}. CEUR-WS. https://ids-pub.bsz-bw.de/frontdoor/index/index/docId/11561.

Bubenhofer, Noah. 2006--2024. Einführung in die Korpuslinguistik: Praktische Grundlagen und Werkzeuge. http://www.bubenhofer.com/korpuslinguistik/.

Bubenhofer, Noah. 2009. \emph{Sprachgebrauchsmuster: Korpuslinguistik als Methode der Diskurs- und Kulturanalyse} (Sprache Und Wissen). Vol. 4. Berlin: De Gruyter. (29 April, 2019).

Bubenhofer, Noah. 2013. Quantitativ informierte qualitative Diskursanalyse: Korpuslinguistische Zugänge zu Einzeltexten und Serien. In Kersten Sven Roth \& Carmen Spiegel (eds.), \emph{Angewandte Diskurslinguistik} (Diskursmuster), 109--134. Berlin: Akademie.

Bubenhofer, Noah. 2016. Drei Thesen zu Visualisierungspraktiken in den Digital Humanities. \emph{Rechtsgeschichte - Legal History}. Max-Planck-Institut fuer europaeische Rechtsgeschichte, Klostermann Verlag. http://dx.doi.org/10.12946/rg24/351-355.

Bubenhofer, Noah. 2017. 4. Kollokationen, n-Gramme, Mehrworteinheiten. In Kersten Sven Roth, Martin Wengeler \& Alexander Ziem (eds.), \emph{Handbuch Sprache in Politik und Gesellschaft}. Berlin, Boston: De Gruyter. https://doi.org/10.1515/9783110296310-004.

Bubenhofer, Noah. 2018a. Wenn „Linguistik`` in „Korpuslinguistik`` bedeutungslos wird. In Joachim Gessinger, Angelika Redder, Ulrich Schmitz \& Wilfried Stölting (eds.), \emph{Korpuslinguistik} (OBST - Osnabrücker Beiträge zur Sprachtheorie 92), 17--30. Duisburg: Universitätsverlag Rhein-Ruhr.

Bubenhofer, Noah. 2018b. 9. Diskurslinguistik und Korpora. In Ingo H. Warnke (ed.), \emph{Handbuch Diskurs}, 208--241. Berlin, Boston: De Gruyter.

Bubenhofer, Noah. 2020a. \emph{Visuelle Linguistik: Zur Genese, Funktion und Kategorisierung von Diagrammen in der Sprachwissenschaft}. De Gruyter. https://doi.org/10.1515/9783110698732.

Bubenhofer, Noah. 2020b. Semantische Äquivalenz in Geburtserzählungen: Anwendung von Word Embeddings. \emph{Zeitschrift für germanistische Linguistik} 48(3). 562--589. https://doi.org/10.1515/zgl-2020-2014.

Bubenhofer, Noah. 2022. Exploration semantischer Räume im Corona-Diskurs. In Heidrun Kämper \& Albrecht Plewnia (eds.), \emph{Sprache in Politik und Gesellschaft}, 197--216. De Gruyter. https://doi.org/10.1515/9783110774306-013.

Bubenhofer, Noah, Selena Calleri \& Philipp Dreesen. 2019. Politisierung in rechtspopulistischen Medien: Wortschatzanalyse und Word Embeddings. \emph{Osnabrücker Beiträge zur Sprachtheorie}. Verein zur Förderung der Sprachwissenschaft in Forschung und Ausbildung 95. 211--241.

Bubenhofer, Noah, Daniel Knuchel \& Larissa Schüller. 2022. Digitale Kulturlinguistik: Digitalität als Gegenstand und Methode. \emph{Germanistik in der Schweiz}. https://doi.org/10.24894/1664-2457.00019.

Bubenhofer, Noah, Daniel Knuchel, Livia Sutter, Maaike Kellenberger \& Niclas Bodenmann. 2020. Von Grenzen und Welten: Eine korpuspragmatische COVID-19-Diskursanalyse. \emph{Aptum} 2--3(16). 156--165.

Bubenhofer, Noah \& Joachim Scharloth. 2013. Korpuslinguistische Diskursanalyse: Der Nutzen empirisch-quantitativer Verfahren. In Ulrike Hanna Meinhof, Martin Reisigl \& Ingo Warnke (eds.), \emph{Diskurslinguistik im Spannungsfeld von Deskription und Kritik}, 147--168. Berlin: Akademie. https://doi.org/10.1524/9783050061047.147.

Bubenhofer, Noah \& Joachim Scharloth. 2015. Maschinelle Textanalyse im Zeichen von Big Data und Data-driven Turn -- Überblick und Desiderate. \emph{Zeitschrift für germanistische Linguistik} 43(1). 1--26. https://doi.org/10.1515/zgl-2015-0001.

Burkhardt, Armin. 1998. Deutsche Sprachgeschichte und politische Geschichte. In Werner Besch, Oskar Reichmann \& Stefan Sonderegger (eds.), \emph{Sprachgeschichte: ein Handbuch zur Geschichte der deutschen Sprache und ihrer Erforschung} (Handbücher zur Sprach- und Kommunikationswissenschaft), vol. 1. Halbband, 98--122. Berlin / New York: De Gruyter.

Busse, Dietrich. 2007. Diskurslinguistik als Kontextualisierung: Sprachwissenschaftliche Überlegungen zur Analyse gesellschaftlichen Wissens. In Ingo Warnke (ed.), \emph{Diskurslinguistik nach Foucault: Theorie und Gegenstände} (Linguistik, Impulse \& Tendenzen 25), 81--105. Berlin/New York: De Gruyter.

Busse, Dietrich. 2008. Diskurslinguistik als Epistemologie -- Das verstehensrelevante Wissen als Gegenstand linguistischer Forschung. In Ingo Warnke \& Jürgen Spitzmüller (eds.), \emph{Methoden der Diskurslinguistik: Sprachwissenschaftliche Zugänge zur transtextuellen Ebene}, 57--88. Berlin / New York: De Gruyter. https://doi.org/10.1515/9783110209372.2.57.

Busse, Dietrich. 2012. \emph{Frame-Semantik: ein Kompendium}. Berlin: De Gruyter.

Busse, Dietrich. 2013. Diskurs -- Sprache -- Gesellschaftliches Wissen: Perspektiven einer Diskursanalyse nach Foucault im Rahmen einer Linguistischen Epistemologie. In Dietrich Busse \& Wolfgang Teubert (eds.), \emph{Linguistische Diskursanalyse: neue Perspektiven}, 147--185. Wiesbaden: Springer. https://doi.org/10.1007/978-3-531-18910-9.

Busse, Dietrich. 2015. \emph{Sprachverstehen und Textinterpretation}. Wiesbaden: Springer. https://doi.org/10.1007/978-3-658-07792-1.

Busse, Dietrich. 2017. 3.3 Lexik -- frame-analytisch. In Thomas Niehr, Jörg Kilian \& Martin Wengeler (eds.), \emph{Handbuch Sprache und Politik: in 3 Bänden} (Sprache - Politik - Gesellschaft Band 21.1-21.3), vol. 21.1, 194--220. Bremen: Hempen.

Busse, Dietrich, Michaela Felden \& Detmer Wulf. 2018. \emph{Bedeutungs- und Begriffswissen im Recht: Frame-Analysen von Rechtsbegriffen im Deutschen} (Sprache und Wissen). Berlin / Boston: De Gruyter. https://doi.org/10.1515/9783110574982.

Busse, Dietrich \& Wolfgang Teubert. 1994. Ist Diskurs ein sprachwissenschaftliches Objekt? In Dietrich Busse, Fritz Hermanns \& Wolfgang Teubert (eds.), \emph{Begriffsgeschichte und Diskursgeschichte} (Methodenfragen und Forschungsergebnisse der historischen Semantik), 10--28. Opladen: Westdeutscher Verlag.

Campello, Ricardo J. G. B., Davoud Moulavi \& Joerg Sander. 2013. Density-Based Clustering Based on Hierarchical Density Estimates. In Jian Pei, Vincent S. Tseng, Longbing Cao, Hiroshi Motoda \& Guandong Xu (eds.), \emph{Advances in Knowledge Discovery and Data Mining} (Lecture Notes in Computer Science), 160--172. Berlin, Heidelberg: Springer. https://doi.org/10.1007/978-3-642-37456-2\_14.

Charniak, Eugene. 1972. \emph{Toward a model of children's story comprehension.} Massachusetts Institute of Technology Thesis. https://dspace.mit.edu/handle/1721.1/13796.

Clarke, Isobelle, Gavin Brookes \& Tony McEnery. 2022. Keywords through Time: Tracking Changes in Press Discourses of Islam. \emph{International Journal of Corpus Linguistics}. https://doi.org/10.1075/ijcl.22011.cla.

Clarke, Isobelle, Tony McEnery \& Gavin Brookes. 2021. Multiple Correspondence Analysis, Newspaper Discourse and Subregister: A Case Study of Discourses of Islam in the British Press. \emph{Register Studies}. John Benjamins 3(1). 144--171. https://doi.org/10.1075/rs.20024.cla.

Clear, Jeremy. 1993. From Firth Principles --- Computational Tools for the Study of Collocation. In Mona Baker, Gill Francis \& Elena Tognini-Bonelli (eds.), \emph{Text and Technology: In honour of John Sinclair}, 271--292. John Benjamins Publishing Company. https://doi.org/10.1075/z.64.18cle.

Czulo, Oliver, Dominic Nyhuis \& Adam Weyell. 2020. Der Einfluss extremistischer Gewaltereignisse auf das Framing von Extremismen auf SPIEGEL Online. \emph{Journal für Medienlinguistik}. Journal für Medienlinguistik 3(1). 14--45. https://doi.org/10.21248/JFML.2020.11.2.

Devlin, Jacob, Ming-Wei Chang, Kenton Lee \& Kristina Toutanova. 2019. BERT: Pre-training of Deep Bidirectional Transformers for Language Understanding. In \emph{Proceedings of the 2019 Conference of the North American Chapter of the Association for Computational Linguistics: Human Language Technologies, Volume 1 (Long and Short Papers)}, 4171--4186. Minneapolis, Minnesota: Association for Computational Linguistics. https://doi.org/10.18653/v1/N19-1423.

Di Carlo, Valerio, Federico Bianchi \& Matteo Palmonari. 2019. Training Temporal Word Embeddings with a Compass. \emph{Proceedings of the AAAI Conference on Artificial Intelligence} 33. 6326--6334. https://doi.org/10.1609/aaai.v33i01.33016326.

Dijk, Teun A. van. 2008. \emph{Discourse and Context: A Sociocognitive Approach}. Cambridge: Cambridge University Press. https://doi.org/10.1017/CBO9780511481499.

Dölemeyer, Anne \& Anne Mehrer. 2011. Einleitung: Ordnung.Macht.Extremismus. In Forum für Kritische Rechtsextremismusforschung (ed.), \emph{Ordnung. Macht. Extremismus: Effekte und Alternativen des Extremismus-Modells}, 7--32. 1. Auflage. Wiesbaden: VS Verlag.

Dreesen, Philipp, Julia Krasselt, Maren Runte \& Peter Stücheli-Herlach. 2023. Operationalisierung der diskurslinguistischen Kategorie ,Akteur`. In Matthias Meiler \& Martin Siefkes (eds.), \emph{Linguistische Methodenreflexion im Aufbruch}, 263--294. De Gruyter. https://doi.org/10.1515/9783111043616-010.

Drommler, Michael. 2024. \emph{Nationale Identität in der Berliner Republik 1998--2007: Ein framesemantischer Zugang zum Wissen gesellschaftlicher Selbstverständigung}. De Gruyter. https://doi.org/10.1515/9783110566246.

Dubossarsky, Haim, Y. Tsvetkov, C. Dyer \& Eitan Grossman. 2015. A Bottom Up Approach to Category Mapping and Meaning Change. In \emph{Proceedings of the NetWordS Final Conference}, vol. 1347, 66--70.

Dunning, Ted. 1993. Accurate Methods for the Statistics of Surprise and Coincidence. (Ed.) Julia Hirschberg. \emph{Computational Linguistics}. Cambridge, MA: MIT Press 19(1). 61--74.

Dutra, Lívia, Arthur Lorenzi, Lorena Larré, Frederico Belcavello, Ely Matos, Amanda Pestana, Kenneth Brown, et al. 2023. Building a Frame-Semantic Model of the Healthcare Domain: Towards the identification of gender-based violence in public health data. In \emph{Anais do XIV Simpósio Brasileiro de Tecnologia da Informação e da Linguagem Humana (STIL 2023)}, 338--346. Brasil: Sociedade Brasileira de Computação. https://doi.org/10.5753/stil.2023.234145.

Ebling, Sahra, Joachim Scharloth, Tobias Dussa \& Noah Bubenhofer. 2013. Gibt es eine Sprache des politischen Extremismus? In Frank Liedtke (ed.), \emph{Die da oben: Texte, Medien, Partizipation} (Sprache - Politik - Gesellschaft), vol. 10, 43--67. Bremen: Hempen.

Egbert, Jesse, Tove Larsson \& Douglas Biber. 2020. \emph{Doing Linguistics with a Corpus: Methodological Considerations for the Everyday User}. 1st edn. Cambridge University Press. https://doi.org/10.1017/9781108888790.

Ehrmanntraut, Anton, Thora Hagen, Leonard Konle \& Fotis Jannidis. 2021. Type- and Token-based Word Embeddings in the Digital Humanities. In \emph{Proceedings of the Conference on Computational Humanities Research 2021}. https://ceur-ws.org/Vol-2989/long\_paper35.pdf.

Entman, Robert M. 1993. Framing: Toward Clarification of a Fractured Paradigm. \emph{Journal of Communication} 43(4). 51--58. https://doi.org/10.1111/j.1460-2466.1993.tb01304.x.

Evert, Stefan. 2005. \emph{The Statistics of Word Cooccurrences: Word Pairs and Collocations}. Stuttgart: Universität Stuttgart. http://dx.doi.org/10.18419/opus-2556.

Evert, Stefan. 2009. 58. Corpora and collocations. In \emph{58. Corpora and collocations}, 1212--1248. De Gruyter Mouton. https://doi.org/10.1515/9783110213881.2.1212.

Evert, Stefan \& Andrew Hardie. 2011. Twenty-first century Corpus Workbench: Updating a query architecture for the new millennium. In \emph{Proceedings of the Corpus Linguistics 2011 conference}. University of Birmingham, UK.

Fankhauser, Peter \& Marc Kupietz. 2019. Analyzing domain specific word embeddings for a large corpus of contemporary German. Application/pdf. In \emph{International Corpus Linguistics Conference, Cardiff, Wales, UK, July 22-26, 2019}. Cardiff: Mannheim. https://doi.org/10.14618/IDS-PUB-9117.

Fankhauser, Peter \& Marc Kupietz. 2022. Count-Based and Predictive Language Models for Exploring DeReKo. In Piotr Bański, Adrien Barbaresi, Simon Clematide, Marc Kupietz \& Harald Lüngen (eds.), \emph{Proceedings of the LREC 2022 Workshop on Challenges in the Management of Large Corpora}, 27--31. Paris: European Language Resources Association (ELRA).

Fasold, Ralph W. 1990. \emph{The Sociolinguistics of Language}. Oxford: Blackwell.

Feiks, Markus. 2019. \emph{Empirische Sozialforschung mit Python: Daten automatisiert sammeln, auswerten, aufbereiten}. Wiesbaden: Springer Fachmedien Wiesbaden. https://doi.org/10.1007/978-3-658-25877-1.

Feilke, Helmuth. 1993. Sprachlicher Common sense und Kommunikation: Über den ``gesunden Menschenverstand'', die Prägung der Kompetenz und die idiomatische Ordnung des Verstehens. \emph{Der Deutschunterricht} 6. 6--21. https://doi.org/10.22029/jlupub-16132.

Feilke, Helmuth. 1996. \emph{Sprache als soziale Gestalt: Ausdruck, Prägung und die Ordnung der sprachlichen Typik}. 1. Aufl. Frankfurt am Main: Suhrkamp.

Feilke, Helmuth. 2003. Textroutine, Textsemantik und sprachliches Wissen. In \emph{Sprache und mehr. Ansichten einer Linguistik der sprachlichen Praxis} (Germanistische Linguistik 245), 209--229. Tübingen: Niemeyer.

Feilke, Helmuth \& Angelika Linke. 2009. Oberfläche und Performanz: Zur Einleitung. In Angelika Linke \& Helmuth Feilke (eds.), \emph{Oberfläche und Performanz: Untersuchungen zur Sprache als dynamischer Gestalt} (Reihe Germanistische Linguistik 283), 3--18. Tübingen: Niemeyer.

Felder, Ekkehard, Marcus Müller \& Friedemann Vogel. 2012. Korpuspragmatik: Paradigma zwischen Handlung, Gesellschaft und Kognition. In Ekkehard Felder, Marcus Müller \& Friedemann Vogel (eds.), \emph{Korpuspragmatik: thematische Korpora als Basis diskurslinguistischer Analysen} (Linguistik, Impulse \& Tendenzen 44), 3--30. Berlin\,/ Boston: De Gruyter.

Fillmore, Charles J. 1965. Entailment Rules in Semantic Theory. \emph{The Ohio State University Research Foundation Project on Linguistic Analysis} Report No. 10. 60--82.

Fillmore, Charles J. 1968. The Case for Case. In Emmon Bach \& Robert Thomas Harms (eds.), \emph{Universals in linguistic theory}, 1--88. New York, NY: Holt, Rinehart \& Winston.

Fillmore, Charles J. 1975. An Alternative to Checklist Theories of Meaning. In \emph{Annual Meeting of the Berkeley Linguistics Society}, 123--131. https://doi.org/10.3765/bls.v1i0.2315.

Fillmore, Charles J. 1976. Frame Semantics and the Nature of Language. In Steven R. Harnad, Horst D. Steklis \& Jane Lancaster (eds.), \emph{Origins and Evolution of Language and Speech} (Annals of the New York Academy of Sciences 280), 20--32. New York: Linguistic Society of America. https://doi.org/10.1111/j.1749-6632.1976.tb25467.x.

Fillmore, Charles J. 1982. Frame Semantics. In The linguistic Society of Korea (ed.), \emph{Linguistics in the Morning Calm: Selected Papers from SICOL-1981}, 111--137. Seoul: Hanshin.

Fillmore, Charles J. 1985. Frames and the Semantics of Understanding. \emph{Quaderni di Semantica} 6(2). 222--254.

Fillmore, Charles J. \& Collin Baker. 2010. A Frames Approach to Semantic Analysis. In Bernd Heine \& Heiko Narrog (eds.), \emph{The Oxford Handbook of Linguistic Analysis} (Oxford Handbooks in Linguistics), 313--339. Oxford: Oxford University Press.

Fillmore, Charles J., Miriam R. L. Petruck, Josef Ruppenhofer \& Abby Wright. 2003. Framenet in Action: The Case of Attaching. \emph{International Journal of Lexicography} 16(3). 297--332. https://doi.org/10.1093/ijl/16.3.297.

Firth, John Rupert. 1957a. A Synopsis of Linguistic Theory 1930--1955. In \emph{Studies in linguistic analysis} (Special Volume of the Philological Society), 1--32. Oxford: Blackwell.

Firth, John Rupert. 1957b. \emph{Papers in Linguistics: 1934-1951}. London: Oxford University Press.

Firth, John Rupert. 1968. \emph{Selected Papers of J. R. Firth 1952-59} (Longmans' Linguistics Library). (Ed.) Frank R. Palmer. Longmans.

Fraas, Claudia, Stefan Meier \& Christian Pentzold. 2010. Konvergenz an den Schnittstellen unterschiedlicher Kommunikationsformen: Ein Frame-basierter analytischer Zugriff. In Hans-Jürgen Bucher, Thomas Gloning \& Katrin Lehnen (eds.), \emph{Neue Medien -- Neue Formate: Ausdifferenzierung und Konvergenz in der Medienkommunikation}, 227--256. Frankfurt am Main/New York: Campus.

Gabrielatos, Costas. 2018. Keyness Analysis: Nature, Metrics and Techniques. In Charlotte Taylor \& Anna Marchi (eds.), \emph{Corpus Approaches to Discourse}, 225--258. Oxford: Routledge.

Gardt, Andreas. 2007. Diskursanalyse: Aktueller theoretischer Ort und methodische Möglichkeiten. In Ingo Warnke (ed.), \emph{Diskurslinguistik nach Foucault: Theorie und Gegenstände} (Linguistik, Impulse \& Tendenzen 25), 27--51. Berlin/New York: De Gruyter.

Gardt, Andreas. 2013. Textanalyse als Basis der Diskursanalyse: Theorie und Methoden. In \emph{Textanalyse als Basis der Diskursanalyse: Theorie und Methoden}, 29--56. De Gruyter. https://doi.org/10.1515/9783110289954.29.

Gillings, Mathew \& Andrew Hardie. 2023. The Interpretation of Topic Models for Scholarly Analysis: An Evaluation and Critique of Current Practice. \emph{Digital Scholarship in the Humanities} 38(2). 530--543. https://doi.org/10.1093/llc/fqac075.

Gillings, Mathew, Mark Learmonth \& Gerlinde Mautner. 2024. Taking the Road Less Travelled: How Corpus‐Assisted Discourse Studies Can Enrich Qualitative Explorations of Large Textual Datasets. \emph{British Journal of Management} 1467-8551.12816. https://doi.org/10.1111/1467-8551.12816.

Gillings, Mathew, Gerlinde Mautner \& Paul Baker. 2023. \emph{Corpus-Assisted Discourse Studies}. 1st edn. Cambridge University Press. https://doi.org/10.1017/9781009168144.

Goffman, Erving. 1974. \emph{Frame Analysis: An Essay on the Organization of Experience} (Harper Colophon Books). New York: Harper \& Row.

Grand, Gabriel, Idan Asher Blank, Francisco Pereira \& Evelina Fedorenko. 2022. Semantic Projection Recovers Rich Human Knowledge of Multiple Object Features from Word Embeddings. \emph{Nature Human Behaviour}. https://doi.org/10.1038/s41562-022-01316-8.

Granger, Sylviane. 2018. Formulaic Sequences in Learner Corpora: Collocations and Lexical Bundles. In Anna Siyanova-Chanturia \& Ana Pellicer-Sánchez (eds.), \emph{Understanding Formulaic Language: A Second Language Acquisition Perspective}, 228--247. Routledge. https://doi.org/10.4324/9781315206615-13.

Harris, Zellig S. 1954. Distributional Structure. \emph{WORD} 10(2--3). 146--162. https://doi.org/10.1080/00437956.1954.11659520.

Helbig, Gerhard \& Joachim Buscha. 2013. \emph{Deutsche Grammatik: Ein Handbuch für den Ausländerunterricht}. Stuttgart: Klett Sprachen GmbH.

Hirscher, Gerhard \& Eckhard Jesse (eds.). 2013. Extremismus in Deutschland. In \emph{Extremismus in Deutschland: Schwerpunkte, Vergleiche, Perspektiven} (Extremismus Und Demokratie Bd. 26), 9--24. 1. Aufl. Baden-Baden: Nomos.

Hunston, Susan \& Gill Francis. 2000. \emph{Pattern Grammar: A Corpus-Driven Approach to the Lexical Grammar of English} (Studies in Corpus Linguistics 4). Amsterdam: Benjamins.

Jacomy, Mathieu, Tommaso Venturini, Sebastien Heymann \& Mathieu Bastian. 2014. ForceAtlas2, a Continuous Graph Layout Algorithm for Handy Network Visualization Designed for the Gephi Software. \emph{PLOS ONE}. Public Library of Science 9(6). e98679. https://doi.org/10.1371/journal.pone.0098679.

Jäger, Siegfried. 1999. Einen Königsweg gibt es nicht. In Hannelore Bublitz, Andrea D. Bührmann, Christine Hanke \& Andrea Seier (eds.), \emph{Das Wuchern der Diskurse: Perspektiven der Diskursanalyse Foucaults}, 136--147. Frankfurt am Main: Campus.

Jäger, Siegfried. 2012. \emph{Kritische Diskursanalyse: Eine Einführung}. Münster: Unrast.

Jarren, Otfried \& Martina Vogel. 2011. „Leitmedien`` als Qualitätsmedien. Theoretisches Konzept und Indikatoren. In Roger Blum, Heinz Bonfadelli, Kurt Imhof \& Otfried Jarren (eds.), \emph{Krise der Leuchttürme öffentlicher Kommunikation: Vergangenheit und Zukunft der Qualitätsmedien}, 17--29. Wiesbaden: VS Verlag für Sozialwissenschaften. https://doi.org/10.1007/978-3-531-93084-8\_2.

Jaschke, Hans-Gerd. 2001. \emph{Rechtsextremismus und Fremdenfeindlichkeit}. Wiesbaden: VS Verlag für Sozialwissenschaften. https://doi.org/10.1007/978-3-322-80839-4.

Jesse, Eckhard. 2006. Grenzen des Demokratieschutzes in der offenen Gesellschaft -- Das Gebot der Äquidistanz gegenüber politischen Extremismen. In Uwe Backes \& Eckhard Jesse (eds.), \emph{Gefährdungen der Freiheit} (Schriften Des Hannah-Arendt-Instituts Für Totalitarismusforschung), vol. 29, 493--520. 1st edn. Göttingen: Vandenhoeck \& Ruprecht. https://doi.org/10.13109/9783666369056.493.

Jesse, Eckhard. 2018. Grundlagen. In Eckhard Jesse \& Tom Mannewitz (eds.), \emph{Extremismusforschung: Handbuch für Wissenschaft und Praxis}, 23--58. 1. Auflage. Baden-Baden: Nomos.

Kailitz, Steffen. 2004. \emph{Politischer Extremismus in der Bundesrepublik Deutschland: eine Einführung} (Lehrbuch). 1. Auflage. Wiesbaden: VS Verlag für Sozialwissenschaften.

Kalwa, Nina. 2013. \emph{Das Konzept\,»Islam«: Eine diskurslinguistische Untersuchung}. Walter de Gruyter.

Kalwa, Nina. 2019. Die Konstitution von Konzepten in Diskursen: Zoom als Methode der diskurslinguistischen Bedeutungsanalyse. In Jürgen Schiewe, Thomas Niehr \& Sandro M. Moraldo (eds.), \emph{Sprach(kritik)kompetenz als Mittel demokratischer Willensbildung: sprachliche In- und Exklusionsstrategien als gesellschaftliche Herausforderung} (Greifswalder Beiträge zur Linguistik Band 12), 11--25. Bremen: Hempen Verlag.

Kämper, Heidrun \& Britt-Marie Schuster (eds.). 2022. \emph{Im Nationalsozialismus: Praktiken -- Kommunikation -- Diskurse. Teil 1}. 1st edn. Göttingen: V\&R unipress. https://doi.org/10.14220/9783737013475.

Katz, Jerrold J. \& Jerry A. Fodor. 1963. The Structure of a Semantic Theory. \emph{Language}. Linguistic Society of America 39(2). 170--210. https://doi.org/10.2307/411200.

Keller, Rudi. 2003. \emph{Sprachwandel: von der unsichtbaren Hand in der Sprache} (UTB Linguistik 1567). Tübingen / Basel: Francke.

Kiess, Johannes. 2011. Rechtsextrem -- extremistisch -- demokratisch? Der prekäre Begriff\,»Rechtsextremismus«\,in der Einstellungsforschung. In Forum für Kritische Rechtsextremismusforschung (ed.), \emph{Ordnung. Macht. Extremismus: Effekte und Alternativen des Extremismus-Modells}, 240--260. 1. Auflage. Wiesbaden: VS Verlag.

Kilgarriff, Adam, Vít Baisa, Jan Bušta, Miloš Jakubíček, Vojtěch Kovář, Jan Michelfeit, Pavel Rychlý \& Vít Suchomel. 2014. The Sketch Engine: Ten Years On. \emph{Lexicography} 1(1). 7--36. https://doi.org/10.1007/s40607-014-0009-9.

Klein, Josef. 1991. Kann man Begriffe besetzen? In Frank Liedtke, Martin Wengeler, Karin Böke, \& Arbeitsgemeinschaft Sprache in der Politik (eds.), \emph{Begriffe besetzen: Strategien des Sprachgebrauchs in der Politik}, 44--69. Opladen: Westdeutscher Verlag.

Klein, Josef. 2018. Frame und Framing: Frametheoretische Konsequenzen aus Praxis und Analyse strategischen politischen Framings. In Alexander Ziem, Lars Inderelst \& Detmer Wulf (eds.), \emph{Frames interdisziplinär: Modelle, Anwendungsfelder, Methoden.} (Proceedings in Language and Cognition), vol. 2, 289--330. Düsseldorf: Düsseldorf University Press.

Klemperer, Victor. 1947. \emph{LTI: Notizbuch eines Philologen}. Berlin: Aufbau.

Knuchel, Daniel \& Noah Bubenhofer. 2023. Machine Learning und Korpuspragmatik. Word Embeddings als Beispiel für einen kreativen Umgang mit NLP-Tools. In Simon Meier-Vieracker, Lars Bülow, Konstanze Marx \& Robert Mroczynski (eds.), \emph{Digitale Pragmatik} (Digitale Linguistik), 213--235. Berlin, Heidelberg: Springer. https://doi.org/10.1007/978-3-662-65373-9\_10.

Kohonen, Teuvo. 2001. \emph{Self-Organizing Maps} (Springer Series in Information Sciences). Vol. 30. Berlin, Heidelberg: Springer Berlin Heidelberg. https://doi.org/10.1007/978-3-642-56927-2.

Konerding, Klaus-Peter. 1993. \emph{Frames und lexikalisches Bedeutungswissen: Untersuchungen zur linguistischen Grundlegung einer Frametheorie und zu ihrer Anwendung in der Lexikographie} (Reihe Germanistische Linguistik). Tübingen: Zugl.: Heidelberg, Univ., Diss., 1992.

Kozlowski, Austin C., Matt Taddy \& James A. Evans. 2019. The Geometry of Culture: Analyzing the Meanings of Class through Word Embeddings. \emph{American Sociological Review} 84(5). 905--949. https://doi.org/10.1177/0003122419877135.

Kupietz, Marc \& Holger Keibel. 2009. Gebrauchsbasierte Grammatik: Statistische Regelhaftigkeit. In Marek Konopka \& Bruno Strecker (eds.), \emph{Deutsche Grammatik - Regeln, Normen, Sprachgebrauch}, 33--50. Berlin, New York: Walter de Gruyter. https://doi.org/10.1515/9783110217360.1.33.

Kupietz, Marc, Harald Lüngen, Paweł Kamocki \& Andreas Witt. 2018. The German Reference Corpus DeReKo: New Developments -- New Opportunities. In Nicoletta Calzolari, Khalid Choukri, Christopher Cieri, Thierry Declerck, Sara Goggi, Koiti Hasida, Hitoshi Isahara, et al. (eds.), \emph{Proceedings of the eleventh international conference on language resources and evaluation (LREC 2018), 7-12 May 2018, Miyazaki, Japan}, 8. Paris: European language resources association (ELRA).

Lakoff, George. 2004. \emph{Don't Think of an Elephant!: Know your Values and Frame the Debate}. Chelsea Green.

Lasch, Alexander. 2014. Zur Vereinbarkeit von diskurslinguistisch motivierter Sprachgeschichtsschreibung und maschineller Sprachanalyse am Beispiel des „Islamismus``-Diskurses. In Vilmos Ágel \& Andreas Gardt (eds.), \emph{Paradigmen der aktuellen Sprachgeschichtsforschung} (Jahrbuch für germanistische Sprachgeschichte 5), 231--249. Berlin, Boston: de Gruyter.

Lemnitzer, Lothar \& Heike Zinsmeister. 2015. \emph{Korpuslinguistik: Eine Einführung} (Narr Studienbücher). Tübingen: Narr.

Lenci, Alessandro. 2018. Distributional Models of Word Meaning. \emph{Annual Review of Linguistics} 4(1). 151--171. https://doi.org/10.1146/annurev-linguistics-030514-125254.

Lenci, Alessandro, Magnus Sahlgren, Patrick Jeuniaux, Amaru Cuba Gyllensten \& Martina Miliani. 2022. A Comparative Evaluation and Analysis of Three Generations of Distributional Semantic Models. \emph{Language Resources and Evaluation} 56(4). 1269--1313. https://doi.org/10.1007/s10579-021-09575-z.

Levy, Omer \& Yoav Goldberg. 2014. Neural word embedding as implicit matrix factorization. In \emph{Proceedings of the 27th International Conference on Neural Information Processing Systems - Volume 2} (NIPS'14), 2177--2185. Cambridge, MA, USA: MIT Press.

Liedtke, Frank, Karin Böke \& Martin Wengeler (eds.). 1991. \emph{Begriffe besetzen: Strategien des Sprachgebrauchs in der Politik}. Opladen: Westdeutscher Verlag.

Link, Jürgen. 1994. Grenzen des flexiblen Normalismus? In Ernst Schulte-Holtey (ed.), \emph{Grenzmarkierungen: Normalisierung und diskursive Ausgrenzung}, 24--39. Duisburg: DISS.

Linke, Angelika. 2003. Begriffsgeschichte -- Diskursgeschichte -- Sprachgebrauchsgeschichte. In Carsten Dutt (ed.), \emph{Herausforderungen der Begriffsgeschichte} (Beiträge zur Philosophie\,: Neue Folge), 39--49. Heidelberg: Winter.

Lönneker-Rodman, Birte \& Alexander Ziem. 2018. Frames als Repräsentationsformat in modernen Terminologiesystemen. In Alexander Ziem, Lars Inderelst \& Detmer Wulf (eds.), \emph{Frames interdisziplinär: Modelle, Anwendungsfelder, Methoden}, 251--288. De Gruyter. https://doi.org/10.1515/9783110720372-009.

Louw, William Ernest. 1993. Irony in the Text or Insincerity in the Writer? The Diagnostic Potential of Semantic Prosodies. In Mona Baker, Gill Francis \& Elena Tognini-Bonelli (eds.), \emph{Text and Technology: In Honour of John Sinclair}, 157--176. Amsterdam: John Benjamins Publishing Company. https://doi.org/10.1075/z.64.11lou.

MacQueen, James. 1967. Some Methods for Classification and Analysis of Multivariate Observations. In \emph{Proceedings of the Fifth Berkeley Symposium on Mathematical Statistics and Probability}. University of California Press.

Marchi, Anna. 2018. Dividing up the Data: Epistemological, Methodological and Practical Impact of Diachronic Segmentation. In Charlotte Taylor \& Anna Marchi (eds.), \emph{Corpus Approaches to Discourse: A Critical Review}, 174--196. Abingdon New York: Routledge.

Matthes, Jörg. 2009. What's in a Frame? A Content Analysis of Media Framing Studies in the World's Leading Communication Journals, 1990-2005. \emph{Journalism \& Mass Communication Quarterly} 86. 349--367. https://doi.org/10.1177/107769900908600206.

Maurer, Marcus, Simon Kruschinski \& Pablo Jost. 2024. Fehlt da was? Perspektivenvielfalt in den öffentlich-rechtlichen Nachrichtenformaten. Mainz. \url{https://www.polkom.ifp.uni-mainz.de/files/2024/01/pm_perspektivenvielfalt.pdf} {[}29.10.2024{]}.

Meier-Vieracker, Simon. 2024. Racist Discourse in a German Far-Right Blog: A Corpus-Driven Approach Using Word Embeddings. \emph{Discourse \& Society}. SAGE Publications Ltd 35(2). 223--242. https://doi.org/10.1177/09579265231204510.

Meißner, Cordula, Daisy Lange \& Christian Fandrych. 2016. Korpusanalyse. In Daniela Caspari, Friederike Klippel, Michael K. Legutke \& Karen Schramm (eds.), \emph{Forschungsmethoden in der Fremdsprachendidaktik}, 306--319. Tübingen: Narr Francke Attempto.

Mikolov, Tomas, Kai Chen, Greg Corrado \& Jeffrey Dean. 2013. Efficient Estimation of Word Representations in Vector Space. arXiv. https://doi.org/10.48550/ARXIV.1301.3781.

Mikolov, Tomas, Ilya Sutskever, Kai Chen, Greg Corrado \& Jeffrey Dean. 2013. Distributed Representations of Words and Phrases and their Compositionality. arXiv. https://doi.org/10.48550/arXiv.1310.4546.

Minsky, Marvin. 1974. A Framework for Representing Knowledge. In \emph{Artificial Intelligence Memo No. 306}. Massachusetts Institute of Technology, A.I. Laboratory,. http://hdl.handle.net/1721.1/6089.

Moretti, Franco. 2000. Conjectures on World Literature. \emph{New Left Review} (1). 54--68.

Murakami, Akira, Paul Thompson, Susan Hunston \& Dominik Vajn. 2017. `What is this corpus about?': Using topic modelling to explore a specialised corpus. \emph{Corpora}. Edinburgh University Press 12(2). 243--277. https://doi.org/10.3366/cor.2017.0118.

Musi, Elena \& Mark Aakhus. 2019. Framing Fracking: Semantic Frames as Meta-Argumentative Indicators for Knowledge-Driven Argument Mining of Controversies. \emph{Journal of Argumentation in Context}. John Benjamins Publishing Company 8(1). 112--135. https://doi.org/10.1075/jaic.18016.mus.

Niehr, Thomas. 2019. Populismus - der Extremismus von heute? \emph{Sprachreport} 35(1). 24--30.

OpenAI. 2023. GPT-4 Technical Report. arXiv. http://arxiv.org/abs/2303.08774.

Oppenhäuser, Holger. 2011. Das Extremismus-Konzept und die Produktion von politischer Normalität. In Forum für kritische Rechtsextremismusforschung (ed.), \emph{Ordnung. Macht. Extremismus}, 35--58. Wiesbaden: VS Verlag für Sozialwissenschaften. https://doi.org/10.1007/978-3-531-93281-1\_2.

Perkuhn, Rainer \& Cyril Belica. 2006. Korpuslinguistik -- Das unbekannte Wesen oder Mythen über Korpora und Korpuslinguistik. \emph{Sprachreport\,: Informationen und Meinungen zur deutschen Sprache}. Institut für Deutsche Sprache 22(1). 2--8.

Pfahl-Traughber, Armin. 2017. Der Erkenntnisgewinn der vergleichenden Extremismusforschung: Besonderheiten, Entwicklungen, Prognosen und Ursachen. In Ralf Altenhof, Sarah Bunk \& Melanie Piepenschneider (eds.), \emph{Politischer Extremismus im Vergleich} (Schriftenreihe Politische Bildung der Konrad-Adenauer-Stiftung e.V.), vol. 3, 45--73. LIT Verlag Münster.

Phillips, Martin. 1989. \emph{Lexical Structure of Text}. English Language Research.

Polenz, Peter von. 1985. \emph{Deutsche Satzsemantik: Grundbegriffe des Zwischen-den-Zeilen-Lesens}. Berlin / New York: De Gruyter.

Rayson, Paul. 2008. From Key Words to Key Semantic Domains. \emph{International Journal of Corpus Linguistics}. John Benjamins 13(4). 519--549. https://doi.org/10.1075/ijcl.13.4.06ray.

Řehůřek, Radim \& Petr Sojka. 2010. Software Framework for Topic Modelling with Large Corpora. In \emph{Proceedings of the LREC 2010 Workshop on New Challenges for NLP Frameworks}, 45--50. Valletta, Malta: ELRA.

Reichardt, Joerg \& Stefan Bornholdt. 2006. Statistical Mechanics of Community Detection. \emph{Physical Review E} 74(1). https://doi.org/10.1103/PhysRevE.74.016110.

Rodatz, Matthias \& Jana Scheuring. 2011. ›Integration als Extremismusprävention‹: Rassistische Effekte der ›wehrhaften Demokratie‹ bei der Konstruktion eines ›islamischen Extremismus‹. In Forum für Kritische Rechtsextremismusforschung (ed.), \emph{Ordnung. Macht. Extremismus: Effekte und Alternativen des Extremismus-Modells}, 163--190. 1. Auflage. Wiesbaden: VS Verlag.

Römer, David \& Martin Wengeler. 2022. Back to the roots! Eine Verteidigungsrede der traditionellen themenbezogenen Diskurslinguistik. \emph{Zeitschrift für Diskursforschung}. Beltz Juventa 10(2). 426--436. https://doi.org/10.3262/ZFD2202426.

Rong, Xin. 2016. word2vec Parameter Learning Explained. arXiv. http://arxiv.org/abs/1411.2738.

Ruppenhofer, Josef, Michael Ellsworth, Miriam R. L. Petruck, Christopher R. Johnson, Collin F. Baker \& Jan Scheffczyk. 2016. \emph{FrameNet II: Extended Theory and Practice}. Revised November 1, 2016. https://framenet2.icsi.berkeley.edu/docs/r1.7/book.pdf.

Sahlgren, Magnus. 2008. The distributional hypothesis. \emph{Rivista die Linguistica} 20(1). 33--53.

Scharloth, Joachim. 2021. \emph{Hässliche Wörter: Hatespeech als Prinzip der neuen Rechten}. Berlin, Heidelberg: Springer Berlin Heidelberg. https://doi.org/10.1007/978-3-662-63502-5.

Scharloth, Joachim. 2022. Faltungen: Die Schließung des rechten Kommunikationssystems aus korpuspragmatischer Perspektive. In \emph{Sprache in Politik und Gesellschaft. Perspektiven und Zugänge}, 217--239. de Gruyter. https://doi.org/10.1515/9783110774306-014.

Scharloth, Joachim, David Eugster \& Noah Bubenhofer. 2013. Das Wuchern der Rhizome: Linguistische Diskursanalyse und Data-driven Turn. In Dietrich Busse \& Wolfgang Teubert (eds.), \emph{Linguistische Diskursanalyse: neue Perspektiven}, 345--380. Wiesbaden: Springer Fachmedien Wiesbaden. https://doi.org/10.1007/978-3-531-18910-9\_11.

Scherer, Carmen. 2006. \emph{Korpuslinguistik}. Heidelberg: Winter.

Schmid, Hans-Jörg. 2003. Collocation: Hard to Pin Down, but Bloody Useful. \emph{Zeitschrift für Anglistik und Amerikanistik} 51(3). 235--258.

Schmid, Helmut. 1994. Probabilistic Part-of-Speech Tagging Using Decision Trees. In \emph{Proceedings of International Conference on New Methods in Language Processing}. Manchester, UK.

Scholz, Ronny \& Alexander Ziem. 2013. Lexikometrie meets FrameNet: das Vokabular der ``Arbeitsmarktkrise'' und der ``Agenda2010'' im Wandel. In Martin Wengeler \& Alexander Ziem (eds.), \emph{Sprachliche Konstruktionen von Krisen: interdisziplinäre Perspektiven auf ein fortwährend aktuelles Phänomen} (Sprache -- Politik -- Gesellschaft 12), 155--184. Bremen: Hempen.

Scholz, Ronny \& Alexander Ziem. 2015. Das Vokabular im diskurshistorischen Vergleich: Skizze einer korpuslinguistischen Untersuchungsheuristik. In Heidrun Kämper \& Ingo H. Warnke (eds.), \emph{Diskurs - interdisziplinär: Zugänge, Gegenstände, Perspektiven}, 281--314. De Gruyter. https://doi.org/10.1515/9783050065281-014.

Schröter, Juliane. 2012. „Wenn Menschen auseinandergehn, So sagen sie: auf Wiedersehn``. Zur soziopragmatischen Geschichte eines Abschiedsgrußes im 19. und 20. Jahrhundert. \emph{Jahrbuch für Germanistische Sprachgeschichte}. De Gruyter 3(1). 361--379. https://doi.org/10.1515/jbgsg-2012-0022.

Schröter, Melani \& Charlotte Taylor (eds.). 2018. \emph{Exploring Silence and Absence in Discourse}. Cham: Springer. https://doi.org/10.1007/978-3-319-64580-3.

Schumacher, Mareike \& Michael Vauth. 2022. Preprocessing mit NLTK. In \emph{forTEXT}. https://fortext.net/routinen/lerneinheiten/preprocessing-mit-nltk. (22 January, 2024).

Sinclair, John. 1991. \emph{Corpus, Concordance, Collocation}. Oxford: Oxford University Press.

Sinclair, John. 1996a. Units of Meaning. \emph{Textus} 9. 75--106.

Sinclair, John. 1996b. SECTION 1: MULTILINGUAL DATABASES: An International Project in Multilingual Lexicography. \emph{International Journal of Lexicography}. Oxford Academic 9(3). 179--196. https://doi.org/10.1093/ijl/9.3.179.

Sinclair, John. 2004. \emph{Trust the Text: Language, Corpus and Discourse}. Annotated Edition. London: Taylor \& Francis.

Sinclair, John. 2005. Meaning in the Framework of Corpus Linguistics. \emph{Lexicographica} 20(2004). 20--32. https://doi.org/10.1515/9783484604674.20.

Spitzmüller, Jürgen \& Ingo Warnke. 2011. \emph{Diskurslinguistik: eine Einführung in Theorien und Methoden der transtextuellen Sprachanalyse} (De Gruyter Studium). Berlin/Boston: De Gruyter.

Stefanowitsch, Anatol. 2020. \emph{Corpus Linguistics: A Guide to the Methodology} (Textbooks in Language Sciences 7). Berlin: Language Science Press. https://doi.org/10.5281/ZENODO.3735822.

Stefanowitsch, Anatol \& Stefan Th Gries. 2003. Collostructions: Investigating the interaction of words and constructions. \emph{International Journal of Corpus Linguistics}. John Benjamins 8(2). 209--243. https://doi.org/10.1075/ijcl.8.2.03ste.

Storjohann, Petra. 2016. \emph{Kookkurrenz aus korpuslinguistischer Sicht} (Literaturhinweise zur Linguistik Band 5). Heidelberg: Universitätsverlag Winter.

Storjohann, Petra \& Melani Schröter. 2013. Präsenz und Absenz lokaler Diskursgebrauchsmuster am Beispiel des deutschen und britischen Krisendiskurses. In Martin Wengeler \& Alexander Ziem (eds.), \emph{Sprachliche Konstruktionen von Krisen: interdisziplinäre Perspektiven auf ein fortwährend aktuelles Phänomen} (Sprache -- Politik -- Gesellschaft 12), 185--208. Bremen: Hempen.

Stubbs, Michael. 1996. \emph{Text and Corpus Analysis: Computer-assisted Studies of Language and Culture} (Language in Society 23). Oxford: Blackwell.

Stulpe, Alexander \& Matthias Lemke. 2016. Blended Reading. In Matthias Lemke \& Gregor Wiedemann (eds.), \emph{Text Mining in den Sozialwissenschaften: Grundlagen und Anwendungen zwischen qualitativer und quantitativer Diskursanalyse}, 17--61. Wiesbaden: Springer Fachmedien. https://doi.org/10.1007/978-3-658-07224-7\_2.

Tesnière, Lucien. 1959. \emph{Élements de syntaxe structurale}. Paris: Klincksieck.

Teubert, Wolfgang. 2013. Die Wirklichkeit des Diskurses. In Dietrich Busse \& Wolfgang Teubert (eds.), \emph{Linguistische Diskursanalyse: neue Perspektiven}, 55--146. Wiesbaden: Springer. https://doi.org/10.1007/978-3-531-18910-9.

Teubert, Wolfgang. 2018. Dietrich Busse und ich. In \emph{Diskurs, Wissen, Sprache}, 31--62. De Gruyter. https://doi.org/10.1515/9783110551853-003.

Teubert, Wolfgang. 2019. Im Kopf oder im Diskurs: wo ist unsere Welt? Sprache und Denken. \emph{tekst i dyskurs - text und diskurs} (12 (2019)). 25--47. https://doi.org/10.7311/tid.12.2019.02.

Thielbörger, Pierre. 2021. Freiheitliche demokratische Grundordnung. In \emph{Handwörterbuch des politischen Systems der Bundesrepublik Deutschland}. https://www.bpb.de/kurz-knapp/lexika/handwoerterbuch-politisches-system/202025/freiheitliche-demokratische-grundordnung/. (27 March, 2024).

Tognini-Bonelli, Elena. 2001. \emph{Corpus Linguistics at Work} (Studies in Corpus Linguistics). Vol. 6. Amsterdam: John Benjamins. https://doi.org/10.1075/scl.6.

Tomasello, Michael. 2006. Konstruktionsgrammatik und früher Erstspracherwerb. In Kerstin Fischer \& Anatol Stefanowitsch (eds.), \emph{Konstruktionsgrammatik: Von der Anwendung zur Theorie} (Konstruktionsgrammatik 1), 19--37. Tübingen: Stauffenburg.

Törnberg, Anton \& Petter Törnberg. 2016. Muslims in Social Media Discourse: Combining Topic Modeling and Critical Discourse Analysis. \emph{Discourse, Context \& Media} 13. 132--142. https://doi.org/10.1016/j.dcm.2016.04.003.

Troiano, Enrica, Roman Klinger \& Sebastian Padó. 2023. On the Relationship between Frames and Emotionality in Text. \emph{Northern European Journal of Language Technology} 9(1). https://doi.org/10.3384/nejlt.2000-1533.2023.4361.

Vachková, Marie \& Cyril Belica. 2008. Self-Organizing Lexical Feature Maps Semiotic Interpretation and Possible Application in Lexicography. \emph{Interdisciplinary Journal for Germanic Linguistics and Semiotic Analysis} 13(2). 223--260.

Warnke, Ingo. 2007a. Diskurslinguistik nach Foucault: Dimensionen einer Sprachwissenschaft jenseits textueller Grenzen. In Ingo Warnke (ed.), \emph{Diskurslinguistik nach Foucault: Theorie und Gegenstände} (Linguistik, Impulse \& Tendenzen 25), 3--24. Berlin/New York: De Gruyter.

Warnke, Ingo (ed.). 2007b. \emph{Diskurslinguistik nach Foucault: Theorie und Gegenstände} (Linguistik, Impulse \& Tendenzen 25). Berlin/New York: De Gruyter.

Wehling, Elisabeth. 2016. \emph{Politisches Framing: Wie eine Nation sich ihr Denken einredet - und daraus Politik macht}. 1st edn. Köln: Herbert von Halem.

Wengeler, Martin. 2022. Warnung vor Framing? Kritische Überlegungen zu Frames und Framing aus polito- und diskurslinguistischer Perspektive. In Kersten Sven Roth \& Martin Wengeler (eds.), \emph{Diesseits und jenseits von Framing: Politikspracheforschung im medialen Diskurs} (Sprache - Politik - Gesellschaft Band 30), 9--29. Hamburg: Buske.

Widdows, Dominic. 2004. \emph{Geometry and Meaning} (Lecture Notes). Stanford: Center for the Study of Language and Information.

Wilk, Nicole M. 2021. Sturm. Nie war mehr Hermeneutik in der Korpuslinguistik (gefragt). \emph{Zeitschrift für Literaturwissenschaft und Linguistik} 51(4). 719--728. https://doi.org/10.1007/s41244-021-00223-1.

Wippermann, Wolfgang. 2009. \emph{Dämonisierung durch Vergleich: DDR und Drittes Reich}. Originalausg., 1. Aufl. Berlin: Rotbuch.

Wodak, Ruth \& Martin Reisigl. 2015. Discourse and Racism. In \emph{The Handbook of Discourse Analysis}, 576--596. John Wiley \& Sons, Ltd. https://doi.org/10.1002/9781118584194.ch27.

Yamada, Kosuke, Ryohei Sasano \& Koichi Takeda. 2021. Semantic Frame Induction using Masked Word Embeddings and Two-Step Clustering. In Chengqing Zong, Fei Xia, Wenjie Li \& Roberto Navigli (eds.), \emph{Proceedings of the 59th Annual Meeting of the Association for Computational Linguistics and the 11th International Joint Conference on Natural Language Processing (Volume 2: Short Papers)}, 811--816. Online: Association for Computational Linguistics. https://doi.org/10.18653/v1/2021.acl-short.102.

Yao, Zijun, Yifan Sun, Weicong Ding, Nikhil Rao \& Hui Xiong. 2018. Dynamic Word Embeddings for Evolving Semantic Discovery. In \emph{Proceedings of the Eleventh ACM International Conference on Web Search and Data Mining - WSDM '18}, 673--681. Marina Del Rey, CA: ACM Press. https://doi.org/10.1145/3159652.3159703.

Zabel, Rebecca. 2015. Sprache, Kultur, Widerstand: Reflexionen über mehr als einen Zusammenhang und wie man Kultur beobachten kann. In Michael Dobstadt, Christian Fandrych \& Ursula Renate Riedner (eds.), \emph{Linguistik und Kulturwissenschaft: zu ihrem Verhältnis aus der Perspektive des Faches Deutsch als Fremd- und Zweitsprache und anderer Disziplinen} (Kulturwissenschaft(En) Als Interdisziplinäres Projekt 9), 65--86. Frankfurt am Main/New York: Peter Lang.

Ziem, Alexander. 2008. \emph{Frames und sprachliches Wissen: kognitive Aspekte der semantischen Kompetenz} (Sprache und Wissen 2). Berlin: de Gruyter.

Ziem, Alexander. 2014. Frames als Prädikations- und Medienrahmen: Auf dem Weg zu einem integrativen Ansatz? In Claudia Fraas, Stefan Meier \& Christian Pentzold (eds.), \emph{Online-Diskurse: Theorien und Methoden transmedialer Online-Diskursforschung.}, 136--172. Köln: Herbert von Halem Verlag.

Ziem, Alexander. 2017. 3. Wortschatz II: quantifizierende Analyseverfahren. In Kersten Sven Roth, Martin Wengeler \& Alexander Ziem (eds.), \emph{Handbuch Sprache in Politik und Gesellschaft}. Berlin, Boston: De Gruyter. https://doi.org/10.1515/9783110296310-003.

Ziem, Alexander. 2018. Der sprachbegabte Mensch ist doch nicht kopflos: Einige Probleme eines radikalen Antikognitivismus. In \emph{Diskurs, Wissen, Sprache}, 63--88. De Gruyter. https://doi.org/10.1515/9783110551853-004.

Ziem, Alexander. 2022. Framing: Genese, Struktur und Problematisierung eines kognitionswissenschaftlichen Konzepts. In Kersten Sven Roth \& Martin Wengeler (eds.), \emph{Diesseits und jenseits von Framing: Politikspracheforschung im medialen Diskurs} (Sprache - Politik - Gesellschaft Band 30), 55--76. Hamburg: Buske.

Ziem, Alexander \& Tim Feldmüller. 2023. Dimensions of Constructional Meanings in the German Constructicon: Why Collo-Profiles Matter. \emph{Yearbook of the German Cognitive Linguistics Association} 11(1). 203--226. https://doi.org/10.1515/gcla-2023-0010.

Ziem, Alexander, Christian Pentzold \& Claudia Fraas. 2018. Medien-Frames als semantische Frames: Aspekte ihrer methodischen und analytischen Verschränkung am Beispiel der ‚Snowden-Affäre`. In Alexander Ziem, Lars Inderelst \& Detmer Wulf (eds.), \emph{Frames interdisziplinär: Modelle, Anwendungsfelder, Methoden} (Proceedings in Language and Cognition 2), 155--182. Düsseldorf: Düsseldorf University Press. https://doi.org/10.1515/9783110720372.

Zimmermann, Jens. 2010. Wissenschaftstheoretische Elemente einer Kritik an der Extremismusforschung und Kritische Diskursanalyse als alternative Perspektive für eine kritische Rechtsextremismusforschung. In Helmut Kellersohn, Martin Dietzsch \& Regina Wamper (eds.), \emph{Rechte Diskurspiraterien: Strategien der Aneignung linker Codes, Symbole und Aktionsformen} (Edition DISS Bd. 28), 264--284. 1. Aufl. Münster: Unrast.

\section{Konkordanzen}\label{konkordanzen}


\begin{table}
K01 taz\_2208663
\begin{tabularx}{\textwidth}{QQQ}
man den Rechten keine Plattform geben darf in den Stuben der Verfassungsschutzämter wird lieber geschreddert & \textbf{Nazis} & sind Soziopathen wer sich ihnen offen entgegenstellt braucht Mut und verdient jede Unterstützung ganz am \\
\end{tabularx}
\end{table}


Angaben gemäß Quelltext der URL.

\section{Abbildungen}\label{abbildungen}

\hyperref[_Ref162016903]{Abbildung 1: Ausschnitt aus dem \FN{Research}-Frame im Berkeley \emph{FrameNet} (Definition und ein Teil der FE). \hyperref[_Ref162016903]{25}}

\hyperref[_Ref162016914]{Abbildung 2: Ergebnis der Abfrage des \FN{Research}-Frames im \emph{FrameGrapher}-Tool in \emph{FrameNet} (\emph{maximum number of generations}: 2, \emph{maximum number of perpipheral children}: 5). \hyperref[_Ref162016914]{26}}

\hyperref[_Ref162016924]{Abbildung 3: Abbildung 7-30 aus Busse (2012: 744): Entwurf eines Maus-Type-Frames. \hyperref[_Ref162016924]{32}}

\hyperref[_Ref162017471]{Abbildung 4: Übersicht 5 aus Klein (2018: 312): Merkel-Ausgangs-Frame. \hyperref[_Ref162017471]{47}}

\hyperref[_Ref162016948]{Abbildung 5: Schaubild 7 aus Backes (1989: 252). Das Hufeisen-Modell der Extremismustheorie. \hyperref[_Ref162016948]{67}}

\hyperref[_Ref162016958]{Abbildung 6: Figure 1 aus Yamada et al. (2021: 811). Originale Bildunterschrift: \emph{2D projections of BERT embeddings of verbs (left) and masked verbs (right). Numbers in the figure correspond to numbers in Table 1,} ⬛ \emph{and \textbf{+} are verbs ``get'' and ``acquire'', respectively, and each color indicates ARRIVING, TRANSITION\_TO\_STATE, and GETTING frame.} \hyperref[_Ref162016958]{82}}

\hyperref[_Ref162016965]{Abbildung 7: \emph{Figure} 2c aus Kozlowski et al. (2019: 913): Semantische Projektion der Wortvektoren unterschiedlicher Sportarten auf zwei ‚kulturelle Dimensionen`. \hyperref[_Ref162016965]{88}}

\hyperref[_Ref162016973]{Abbildung 8: \emph{Figure 3} aus Vachková \& Belica (2008: 250--251): Annotierte SOM zu \emph{Tod}. \hyperref[_Ref162016973]{91}}

\hyperref[_Ref180396143]{Abbildung 9: Abbildung 9 aus Maurer et al. (2024: 17). Verortung großer deutscher Medien in Hinblick auf ihre "ideologische Grundpositionierung". \hyperref[_Ref180396143]{102}}

\hyperref[_Ref162016992]{Abbildung 10: Umfang der Subkorpora in Millionen Token. \hyperref[_Ref162016992]{106}}

\hyperref[_Ref162017013]{Abbildung 11: Räumliche Darstellung der \emph{Log-Likelihood}-Werte der Städte im Verhältnis zu zwei Clustern aus Tabelle 4. \hyperref[_Ref162017013]{120}}

\hyperref[_Ref162017046]{Abbildung 12: Abbildung 5.11 aus Biemann et al. (2022: 221): Schematische Darstellung der CBOW-Architektur für die Vorhersage von \emph{als} im Satz \emph{Mozart galt als musikalisches Wunderkind}. \hyperref[_Ref162017046]{124}}

\hyperref[_Ref162017066]{Abbildung 13: Abbildung 6.5 aus Biemann et al. (2022: 270). Illustration des \emph{k-Means}-Algorithmus bei \emph{k} = 2 und 6 Vektoren. \hyperref[_Ref162017066]{129}}

\hyperref[_Ref162016429]{Abbildung 14: \emph{Figure 2a} aus Grand et al. (2022). Korrelation zwischen der Projektion der Word Embeddings unterschiedlicher Tiere auf unterschiedliche semantische Achsen und menschlichen Bewertungen. \hyperref[_Ref162016429]{130}}

\hyperref[_Ref162016479]{Abbildung 15: Nächste paradigmatische Nachbarn zu \emph{extremistische} im Zeitraum 2020 -- 2021 (Dimensionsreduktion mit PCA). \hyperref[_Ref162016479]{132}}

\hyperref[_Ref162016499]{Abbildung 16: Extremismus-Frame des Zeitraumes 2020 -- 2021. \hyperref[_Ref162016499]{141}}

\hyperref[_Ref162016524]{Abbildung 17: Ausschnitt aus dem Politik-Teil-Frame des Zeitraums 2020 -- 2021. \hyperref[_Ref162016524]{141}}

\hyperref[_Ref162016542]{Abbildung 18: Teil des Hanau-Frame des Zeitraums 2020 -- 2021. \hyperref[_Ref162016542]{142}}

\hyperref[_Ref162453186]{Abbildung 19: Rechtsextremismus-Frame des Zeitraums 1999--2001. \hyperref[_Ref162453186]{144}}

\hyperref[_Ref162016812]{Abbildung 20: Tooltip zum Cluster 11\_Uwe\_Mundlos{\breakbar}Uwe\_Böhnhardt{\breakbar}Böhnhardt. \hyperref[_Ref162016812]{161}}

\hyperref[_Ref162439852]{Abbildung 21: Ausschnitt aus dem thematischen Teil-Frame ‚Politik` am Beispiel des Clusters 02\_Nein{\breakbar}Ja. \hyperref[_Ref162439852]{176}}

\hyperref[_Ref162439838]{Abbildung 22: Ausschnitt aus dem thematischen Teil-Frame ‚Krieg \& Islamismus` am Beispiel des Clusters 11\_Militärstützpunkt{\breakbar}Selbstmordattentat{\breakbar}Luftwaffenstützpunkt. \hyperref[_Ref162439838]{178}}

\hyperref[_Ref162439822]{Abbildung 23: Ausschnitt aus dem thematischen Teil-Frame ‚Straftaten \& Juristisches` am Beispiel des Clusters ta\_Geständnis. \hyperref[_Ref162439822]{179}}

\hyperref[_Ref162439808]{Abbildung 24: Ausschnitt aus dem thematischen Teil-Frame ‚Staat, FdGO \& Gesellschaft` am Beispiel des Clusters 20\_Öffentlichkeit{\breakbar}Presse. \hyperref[_Ref162439808]{180}}

\hyperref[_Ref162439766]{Abbildung 25: Ausschnitt aus dem thematischen Teil-Frame ‚Außenpolitik` am Beispiel des Clusters 14\_Reformen. \hyperref[_Ref162439766]{181}}

\hyperref[_Ref162439793]{Abbildung 26: Ausschnitt aus dem thematischen Teil-Frame ‚Demonstrationen \& politischer Extremismus` am Beispiel des Clusters 17\_demonstrantInnen{\breakbar}Demonstrantinnen{\breakbar}Protestierenden. \hyperref[_Ref162439793]{181}}

\hyperref[_Ref162439728]{Abbildung 27: Ausschnitt aus dem thematischen Teil-Frame ‚Flucht \& Migration` am Beispiel des Clusters 14\_Zuwanderung{\breakbar}Einwanderung{\breakbar}Migration. \hyperref[_Ref162439728]{183}}

\hyperref[_Ref162450642]{Abbildung 28: Thematische Verankerung des NPD-Frames. \hyperref[_Ref162450642]{186}}

\hyperref[_Ref162450557]{Abbildung 29: Perspektivierung rechtsextremistischer Akteure durch das FE 99\_Nazis. \hyperref[_Ref162450557]{187}}

\hyperref[_Ref162450533]{Abbildung 30: Perspektivierung rechtsextremistischer Akteure durch das FE 99\_Neonazis{\breakbar}Skinheads{\breakbar}Rechtsextremisten. \hyperref[_Ref162450533]{188}}

\hyperref[_Ref162450504]{Abbildung 31: Perspektivierung des Rechtsextremismus-Frames durch das FE 99\_Fremdenfeindlichkeit{\breakbar}Ausländerfeindlichkeit{\breakbar}Rassismus. \hyperref[_Ref162450504]{189}}

\hyperref[_Ref162441675]{Abbildung 32: Vernetzung des FE 02\_Pds (rechts) im Gesamt-Frame. \hyperref[_Ref162441675]{196}}

\hyperref[_Ref162442308]{Abbildung 33: Neukonfiguration des Framings von Islamismus bzw. Demonstrationen am Beispiel des FE 11\_Demonstranten. \hyperref[_Ref162442308]{216}}

\hyperref[_Ref162457718]{Abbildung 34: Teil der Vernetzung des FE 14\_Silvesternacht im Bereich ‚Straftaten \& Juristisches`. \hyperref[_Ref162457718]{221}}

\hyperref[_Ref162442563]{Abbildung 35: Teil des George-Floyd-Frames. \hyperref[_Ref162442563]{234}}

\section{Tabellen}\label{tabellen}

\hyperref[_Ref162017554]{Tabelle 1: Nächste Nachbarn im Word-Embedding-Modell 2020 -- 2021 zu \emph{Junge}, \emph{Alternative} und \emph{Junge\_Alternative} mit Cosinus-Ähnlichkeit. \hyperref[_Ref162017554]{108}}

\hyperref[_Ref162017306]{Tabelle 2: Type-Anzahlen und Type-Anzahlen mit einer Mindestfrequenz von 25 pro Subkorpus. \hyperref[_Ref162017306]{108}}

\hyperref[_Ref162017574]{Tabelle 3: Beispiel-Output einer Kollokationsanalyse mit \emph{polmineR} zu \emph{Hanau} im Subkorpus 2020 -- 2021. \hyperref[_Ref162017574]{116}}

\hyperref[_Ref162017334]{Tabelle 4: Kollokationsmatrix zu Wort-Clustern des Zeitraumes 2020 -- 2021 mit \emph{Log-Likelihood}-Werten. Dargestellt ist die gerichtete Kollokationsstärke der Zeilen-Cluster zu den Spalten-Clustern. Vermerkte Signifikanzlevel: *p \textless{} .05 (LL \textgreater= 3,84), **p \textless{} .01 (LL \textgreater= 6,63), ***p \textless{} .001 (LL \textgreater= 10,38), ****p \textless{} .0001 (LL \textgreater= 15,13). \hyperref[_Ref162017334]{118}}

\hyperref[_Ref162017640]{Tabelle 5: Nächste Nachbarn einzelner Embeddings mit Cosinus-Ähnlichkeit im zeitlichen Verlauf. \hyperref[_Ref162017640]{127}}

\hyperref[_Ref162017162]{Tabelle 6: Drei Beispiel-Cluster des Zeitraums 2014--2016 und alle enthaltenen Wörter (gereiht nach Nähe zum Cluster-Centroid). \hyperref[_Ref162017162]{134}}

\hyperref[_Ref162017170]{Tabelle 7: Beispiele für Cluster mit unterschiedlicher Distanz zum Extremismus-Centroid. \hyperref[_Ref162017170]{136}}

\hyperref[_Ref162017239]{Tabelle 8: Kollokationen und \emph{next neighbours} im Word-Embedding-Modell zum Cluster 11\_festgenommen. \hyperref[_Ref162017239]{139}}

\hyperref[_Ref162017380]{Tabelle 9: Zu betrachtende Cluster für den medialen Vergleich nach Extremismsuvariante und thematischem Bezug. \hyperref[_Ref162017380]{174}}

\hyperref[_Ref162017395]{Tabelle 10: Thematisch gruppierte Auswahl semantisch äquivalenter Cluster des Zeitraums 2002--2004. \hyperref[_Ref162017395]{194}}

\hyperref[_Ref162017417]{Tabelle 11: Thematisch gruppierte Auswahl semantisch äquivalenter Cluster des Zeitraums 2005--2007. \hyperref[_Ref162017417]{201}}

\hyperref[_Toc162477804]{Tabelle 12: Thematisch gruppierte Auswahl semantisch äquivalenter Cluster des Zeitraums 2008--2010. \hyperref[_Toc162477804]{207}}

\hyperref[_Toc162477805]{Tabelle 13: Thematisch gruppierte Auswahl semantisch äquivalenter Cluster des Zeitraums 2011--2013. \hyperref[_Toc162477805]{214}}

\hyperref[_Toc162477806]{Tabelle 14: Thematisch gruppierte Auswahl semantisch äquivalenter Cluster des Zeitraums 2014--2016. \hyperref[_Toc162477806]{219}}

\hyperref[_Toc162477807]{Tabelle 15: Thematisch gruppierte Auswahl semantisch äquivalenter Cluster des Zeitraums 2017--2019. \hyperref[_Toc162477807]{226}}

\hyperref[_Toc162477808]{Tabelle 16: Thematisch gruppierte Auswahl semantisch äquivalenter Cluster des Zeitraums 2020 -- 2021. \hyperref[_Toc162477808]{231}}

\hyperref[_Ref162017445]{Tabelle 17: Thematisch gruppierte Auswahl semantisch äquivalenter Cluster der \textsc{TAZ}. \hyperref[_Ref162017445]{241}}

\hyperref[_Toc162477810]{Tabelle 18: Thematisch gruppierte Auswahl semantisch äquivalenter Cluster des \textsc{Spiegel}. \hyperref[_Toc162477810]{249}}

\hyperref[_Toc162477811]{Tabelle 19: Thematisch gruppierte Auswahl semantisch äquivalenter Cluster der \textsc{Welt}. \hyperref[_Toc162477811]{254}}